%% Part IV -- Chances (comic pages 28 to 36).
%% Eight arguments for the technology, made properly: the scale argument, the
%% accuracy argument, the bias argument turned around, and then the quiet one
%% on page 36 that a chance also costs something.
%%
%% ONE PANEL CAPTION, ONE SLIDE. The comic sets every topic page as six
%% panels in a 2x3 grid with one caption under each. Each caption takes a
%% slide of its own so the room reads one claim at a time; the six-dot row
%% under it marks which panel is on screen, and six consecutive slides
%% sharing a frame title are one comic page, read in the comic's own order.
%% Captions are verbatim, checked against a render of each page.
%% \setpagethesis opens each run with that page's one line, mine, not theirs.
%% Fragment only: no preamble, \input by ai_for_all.tex.

\sectionsubtitle{pages 28 to 36}
\section{Chances}

% The bottleneck is not storage, it is us. Everything after this page in
% the section follows from panel 3.
\setpagethesis{The bottleneck was never storage. It is us.}
\begin{frame}{Dealing with Big Data}
  \panel{hundreds of us against the volume}
  \panelslide{29}{1}{Our power to analyze the tons of data we are now able to store and to collect can be a bottleneck.}
\end{frame}
\begin{frame}{Dealing with Big Data}
  \panel{hundreds of us against the volume}
  \panelslide{29}{2}{AI helps us deal with large and complex datasets in ways we have never seen before.}
\end{frame}
\begin{frame}{Dealing with Big Data}
  \panel{hundreds of us against the volume}
  \panelslide{29}{3}{Even with hundreds of us analyzing patterns in big data, the sheer volume simply overwhelms our capacities.}
\end{frame}
\begin{frame}{Dealing with Big Data}
  \panel{hundreds of us against the volume}
  \panelslide{29}{4}{Even if we could analyze data, but it proves to be exhausting and tedious, we can use rule-based AI systems -- expert systems -- that do the job for us.}
\end{frame}
\begin{frame}{Dealing with Big Data}
  \panel{hundreds of us against the volume}
  \panelslide{29}{5}{Traditional algorithms are bound to follow the same logic over and over, without flexibility or learning new paths with new data.}
\end{frame}
\begin{frame}{Dealing with Big Data}
  \panel{hundreds of us against the volume}
  \panelslide{29}{6}{Now enter AI. Because AI systems get smarter as more data is given to them, they are well suited for patterns and anomalies over the time.}
\end{frame}

% The lawyer experiment: 20 lawyers, 5 contracts, 30 issues, 4 hours.
% Panel 5 holds all four numbers.
\setpagethesis{Ninety-five against eighty-five. What would you need to know to believe it?}
\begin{frame}{Efficiency}
  \panel{twenty lawyers against one AI}
  \panelslide{30}{1}{In many situations, AI is simply faster and more accurate than we are. Let's look at an example*.}
\end{frame}
\begin{frame}{Efficiency}
  \panel{twenty lawyers against one AI}
  \panelslide{30}{2}{A legal platform carried out a competition between 20 experienced lawyers against a trained AI.}
\end{frame}
\begin{frame}{Efficiency}
  \panel{twenty lawyers against one AI}
  \panelslide{30}{3}{Competitors had to review 5 legal contracts in 4 hours and identify 30 legal issues.}
\end{frame}
\begin{frame}{Efficiency}
  \panel{twenty lawyers against one AI}
  \panelslide{30}{4}{They were rated according to how accurately they identified each issue. Who won?}
\end{frame}
\begin{frame}{Efficiency}
  \panel{twenty lawyers against one AI}
  \panelslide{30}{5}{Humans achieved, on average, an 85\% accuracy rate, the AI 95\%. Humans took 92 minutes to complete the task, the AI\ldots{} 26 seconds.}
\end{frame}
\begin{frame}{Efficiency}
  \panel{twenty lawyers against one AI}
  \panelslide{30}{6}{The human lawyers who competed against the AI in the experiment said, the tasks were very similar to what lawyers do every day.}
\end{frame}

% The sharpest reasoning in the book. Panel 5 is the point: a shared
% distortion does not average out, it compounds.
\setpagethesis{A shared distortion does not average out. It compounds.}
\begin{frame}{Cognitive Biases}
  \panel{gender-bias guiding you through the world}
  \panelslide{31}{1}{Our brain, as a result of evolution, was not optimized for rational decision-making or perfect diagnoses.}
\end{frame}
\begin{frame}{Cognitive Biases}
  \panel{gender-bias guiding you through the world}
  \panelslide{31}{2}{Instead, we strive for a competitive degree of `fitness' in our specific environments; quite fast and often satisfying.}
\end{frame}
\begin{frame}{Cognitive Biases}
  \panel{gender-bias guiding you through the world}
  \panelslide{31}{3}{Our brain power is biologically limited; therefore, we work with heuristics, trust our gut feeling or store \& retrieve wrong evidence.}
\end{frame}
\begin{frame}{Cognitive Biases}
  \panel{gender-bias guiding you through the world}
  \panelslide{31}{4}{We call this cognitive bias. There are many of them and we share them to some degree.}
\end{frame}
\begin{frame}{Cognitive Biases}
  \panel{gender-bias guiding you through the world}
  \panelslide{31}{5}{And this is what makes our distortions have large-scale consequences. Were they random, they would cancel out.}
\end{frame}
\begin{frame}{Cognitive Biases}
  \panel{gender-bias guiding you through the world}
  \panelslide{31}{6}{Used wisely, AI could indeed help us improve our diagnoses and decisions.}
\end{frame}

% Gives the machine the generation and keeps the judgement. The last clause
% of panel 4 is the position, and it has not aged.
\setpagethesis{The machine generates. Judging stays with us, and judging is the part we enjoy.}
\begin{frame}{Creativity}
  \panel{just a generic style transfer}
  \panelslide{32}{1}{AI can save time otherwise wasted on stupid, tedious work for creativity. But it can do even more.}
\end{frame}
\begin{frame}{Creativity}
  \panel{just a generic style transfer}
  \panelslide{32}{2}{When an AI sorts our photos, it is easier to make a photo-calendar with our children (in 28 positions) for our parents as a present.}
\end{frame}
\begin{frame}{Creativity}
  \panel{just a generic style transfer}
  \panelslide{32}{3}{When an AI puts together a playlist of up-and-coming producers of house music, it may inspire our next DJ-set.}
\end{frame}
\begin{frame}{Creativity}
  \panel{just a generic style transfer}
  \panelslide{32}{4}{We can also get creative directly with the AI, compiling data and coding it, in a narrow area. But the AI has no judgement.}
\end{frame}
\begin{frame}{Creativity}
  \panel{just a generic style transfer}
  \panelslide{32}{5}{Judging is an ability we possess. And we love doing it. This is where AI together with us can achieve great things, with a human tweak.}
\end{frame}
\begin{frame}{Creativity}
  \panel{just a generic style transfer}
  \panelslide{32}{6}{The smaller the role of humans in the process of art-creation becomes, the more AI questions traditional concepts of art.}
\end{frame}

% The chance here is not the diagnosis, it is the time the diagnosis frees
% up. Panel 2 is unusually blunt about what that time currently goes on.
\setpagethesis{The chance is not the diagnosis. It is the time the diagnosis frees up.}
\begin{frame}{Connection}
  \panel{the doctor as overpaid bookkeeper}
  \panelslide{33}{1}{Most doctors became doctors because they wanted to connect with people in order to heal them.}
\end{frame}
\begin{frame}{Connection}
  \panel{the doctor as overpaid bookkeeper}
  \panelslide{33}{2}{Today, they often act like overpaid bookkeepers: taking in and spitting back data, prescribing drugs, adjusting doses, ordering tests.}
\end{frame}
\begin{frame}{Connection}
  \panel{the doctor as overpaid bookkeeper}
  \panelslide{33}{3}{And it's still only retrospective, error-prone knowledge.}
\end{frame}
\begin{frame}{Connection}
  \panel{the doctor as overpaid bookkeeper}
  \panelslide{33}{4}{But we cannot connect to doctors who struggle to keep up with messy data from disparate connections. And they cannot connect to us.}
\end{frame}
\begin{frame}{Connection}
  \panel{the doctor as overpaid bookkeeper}
  \panelslide{33}{5}{AI can support proactive, individualized strategies.}
\end{frame}
\begin{frame}{Connection}
  \panel{the doctor as overpaid bookkeeper}
  \panelslide{33}{6}{Our doctors could get to know us better, having the time to learn how a disease uniquely affects us and educate us in the right way.}
\end{frame}

% Robotics without reprogramming: demonstration replaces the rule set,
% which is page 15's move applied to a body.
\setpagethesis{Show the robot instead of programming it: page 15's move, applied to a body.}
\begin{frame}{New Opportunities}
  \panel{the sensor jacket and glove}
  \panelslide{34}{1}{AI allows for fascinating new opportunities. Teaching robots without a single line of coding, for example.}
\end{frame}
\begin{frame}{New Opportunities}
  \panel{the sensor jacket and glove}
  \panelslide{34}{2}{Robots solve complex tasks today -- but these tasks are deterministic. At any point of time, you know in advance where your robot needs to be.}
\end{frame}
\begin{frame}{New Opportunities}
  \panel{the sensor jacket and glove}
  \panelslide{34}{3}{Every time a task or location changes, the robot's software must be re-programmed. This is time-consuming, expensive and difficult.}
\end{frame}
\begin{frame}{New Opportunities}
  \panel{the sensor jacket and glove}
  \panelslide{34}{4}{Even for experienced programmers, it is hard to come up with a set of verbal rules of how to move when movements and settings are tricky or unpredictable.}
\end{frame}
\begin{frame}{New Opportunities}
  \panel{the sensor jacket and glove}
  \panelslide{34}{5}{AI now enables us to teach robots how to move, either together with the robot, or, with jackets and gloves containing sensors.}
\end{frame}
\begin{frame}{New Opportunities}
  \panel{the sensor jacket and glove}
  \panelslide{34}{6}{From the recorded example-movements, automation scripts can be created and optimized -- without us having to write complex code. Fantastic!}
\end{frame}

% The least contested case in the book. Reading, hearing, typing, the home,
% the classroom and the ward, one per panel.
\setpagethesis{Six barriers, one per panel, and not one of them is about convenience.}
\begin{frame}{Inclusion}
  \panel{a wall missing several bricks}
  \panelslide{35}{1}{AI is also about helping us, about improving inclusion and our well-being, by removing barriers.}
\end{frame}
\begin{frame}{Inclusion}
  \panel{a wall missing several bricks}
  \panelslide{35}{2}{AI reads to us, it narrates the world around us and help us `see'. Speech-to-text AI helps us `hear', too.}
\end{frame}
\begin{frame}{Inclusion}
  \panel{a wall missing several bricks}
  \panelslide{35}{3}{AI helps us type (with words) or talks to us. Users of all ages are adopting virtual assistants because it is so simple to learn how to handle them.}
\end{frame}
\begin{frame}{Inclusion}
  \panel{a wall missing several bricks}
  \panelslide{35}{4}{Smart home devices help us do things easily around the house and live independently.}
\end{frame}
\begin{frame}{Inclusion}
  \panel{a wall missing several bricks}
  \panelslide{35}{5}{In the classroom, AI can identify individual challenges and offer more personalized approaches.}
\end{frame}
\begin{frame}{Inclusion}
  \panel{a wall missing several bricks}
  \panelslide{35}{6}{AI virtual nurses or therapists offer 24-hour-support, or someone to talk to in case we are shy.}
\end{frame}

% Closes the section by undercutting it. The chance is convenience and the
% bill is the analogue institution. Filed under Chances, not Risks, on purpose.
\setpagethesis{What did you stop leaving the house for, and what closed because of it?}
\begin{frame}{Comfort}
  \panel{mall\_blueprint.pdf}
  \panelslide{36}{1}{AI satisfies and reinforces our desire for comfort.}
\end{frame}
\begin{frame}{Comfort}
  \panel{mall\_blueprint.pdf}
  \panelslide{36}{2}{Almost all answers to our consumer, communication or information needs are just a mouse click or an app away. And most of the time, AI is involved.}
\end{frame}
\begin{frame}{Comfort}
  \panel{mall\_blueprint.pdf}
  \panelslide{36}{3}{It is tempting not to have to leave our couch at all, no matter what we want to do: eat, work, chat, watch a movie, listen to music\ldots{}}
\end{frame}
\begin{frame}{Comfort}
  \panel{mall\_blueprint.pdf}
  \panelslide{36}{4}{'Pre-AI'-institutions such as shops, local transport, restaurants, libraries, or doctors must react to our new convenience.}
\end{frame}
\begin{frame}{Comfort}
  \panel{mall\_blueprint.pdf}
  \panelslide{36}{5}{All institutions with predominantly analogue offerings are undergoing a tremendous process of change. But it might be worth it.}
\end{frame}
\begin{frame}{Comfort}
  \panel{mall\_blueprint.pdf}
  \panelslide{36}{6}{It is precisely these institutions that are social places where we can come together and exchange ideas. Where we can decelerate. `Waste' time.}
\end{frame}
